\documentclass[11pt]{article}
\usepackage[margin=1.1in]{geometry}
\usepackage{amsmath,amssymb,amsthm}
\usepackage[colorlinks=true,linkcolor=blue,citecolor=blue]{hyperref}

\newtheorem{theorem}{Theorem}[section]
\newtheorem{lemma}[theorem]{Lemma}
\newtheorem{proposition}[theorem]{Proposition}
\newtheorem{corollary}[theorem]{Corollary}
\theoremstyle{definition}
\newtheorem{definition}[theorem]{Definition}
\theoremstyle{remark}
\newtheorem{remark}[theorem]{Remark}

\newcommand{\R}{\mathbb{R}}
\newcommand{\E}{\mathbb{E}}
\newcommand{\PP}{\mathbb{P}}
\newcommand{\emp}{\bar\mu_N}
\newcommand{\vol}{\operatorname{vol}}
\newcommand{\cone}{\operatorname{cone}}
\newcommand{\dist}{\operatorname{dist}}
\newcommand{\nor}{\nu_K}

\title{The interpolation law for the uniform doubling condition\\
in general position}
\author{Yitzchak Shmalo\\
\small Einstein Institute of Mathematics, The Hebrew University of Jerusalem, Jerusalem, Israel\\
\small \texttt{yitzchak.shmalo@gmail.com}}
\date{June 2026}

\begin{document}
\maketitle

\begin{abstract}
The companion paper \cite{companion11} determined the critical scaling of the
uniform doubling condition (UDC) of \cite{ShmaloUDC} for the orthogonal model
$B^n\cap\{x_1\le0,\dots,x_k\le0\}$, where a round ball is cut by mutually
orthogonal half-spaces through the origin.  Here we remove both idealizations.
We treat an admissible configuration: a $C^2$ strictly convex body $K$ with
principal curvatures in a fixed band $[\lambda_{\min},\lambda_{\max}]$,
truncated by $k$ half-spaces $\langle x,a_i\rangle\le c_i$ in general
position, meaning that at each point $p$ of the boundary stratum
$E=\partial K\cap\bigcap_i\{\langle x,a_i\rangle=c_i\}$ the outer normal
$\nor(p)$ and the truncating normals $a_1,\dots,a_k$ are linearly independent
with pairwise angles in $[\theta_0,\pi-\theta_0]$.  The stratum $E$ is then a
$j$-manifold with $j=n-1-k$.  We prove that the critical scaling factor is
\[
\ell_N\;\asymp\;\Bigl(\frac{\log N}{N}\Bigr)^{2/(2n-j)}
\;=\;\Bigl(\frac{\log N}{N}\Bigr)^{2/(n+1+k)},
\]
exponent and logarithm exact, with constants depending only on $n$, $k$,
$\theta_0$ and the curvature band.  The mechanism that replaces the
orthogonal model's axis-parallel windows and its exact reflection fraction
$2^{-k}$ is the \emph{conjugate frame}: the dual basis $f_1,\dots,f_k$ to the
truncating normals inside their span, chosen tangent to $K$ at $p$, in which
the $k$ constraints become the diagonal slab $\{s_i\le0\}$ and the cap
exponent $(n+1+k)/2$ is exposed as angle-independent.  The reflection fraction
is replaced by the normalized solid angle $\omega$ of the normal cone
$N(G,p)=\cone\{\nor(p),a_1,\dots,a_k\}$, which is positive precisely because
the configuration is in general position.  We note that the natural
attempt to reduce to the orthogonal model by a John-position rounding is
provably blocked --- no single linear map can keep $B^n$ spherical and
orthogonalize non-orthogonal normals --- so the conjugate-frame argument is
intrinsic by necessity.  This completes the flat-transversal part of the program
of \cite{companion11,companion6,companion7}.  The curved companion
\cite{companion15curved} proves the same exponent law for each fixed curved
truncation family and isolates the remaining curvature-band-uniform lower/slab
constant problem.
\end{abstract}

\section{Setting and main results}\label{sec:intro}

\subsection{The configuration}

Let $n\ge2$.  Let $K\subset\R^n$ be a compact convex body with $C^2$
boundary whose principal curvatures lie in a fixed band
$0<\lambda_{\min}\le\lambda_{\max}<\infty$; in particular $K$ is strictly
convex, the Gauss map $\nor:\partial K\to S^{n-1}$ is a diffeomorphism, and at
each $p\in\partial K$ the second fundamental form $\mathrm{I\!I}_p$ satisfies
$\lambda_{\min}\,|v|^2\le\mathrm{I\!I}_p(v,v)\le\lambda_{\max}\,|v|^2$ for
tangent $v$.  Let $\mu$ be the normalized Lebesgue measure on $K$.

Fix unit vectors $a_1,\dots,a_k\in S^{n-1}$ and reals $c_1,\dots,c_k$ with
$1\le k\le n-1$ (the case $k=0$ is the smooth body, treated in
\cite{companion3}), and set the \emph{truncation region}
\[
G\;:=\;K\cap\bigcap_{i=1}^{k}\{x:\ \langle x,a_i\rangle\le c_i\}.
\]
We call $a_i$ the truncating normals and $h_k:=\{(a_1,c_1),\dots,(a_k,c_k)\}$
the truncation set.  The boundary of $G$ is stratified by the active
constraint set: for $S\subseteq\{1,\dots,k\}$ the stratum is
$\partial K\cap\bigcap_{i\in S}\{\langle x,a_i\rangle=c_i\}
\cap\bigcap_{i\notin S}\{\langle x,a_i\rangle<c_i\}$ together with the curved
face $\{x\in\mathrm{int}\,K:\,\langle x,a_i\rangle=c_i\ (i\in S)\}$.  The
lowest transversal stratum is
\[
E\;:=\;\partial K\cap\bigcap_{i=1}^{k}\{x:\ \langle x,a_i\rangle=c_i\}.
\]

\begin{definition}[Admissible / general position]\label{def:adm}
The configuration $(K,h_k)$ is \emph{admissible}, or the truncation is in
\emph{general position}, if there is $\theta_0\in(0,\pi/2)$ such that
$E\ne\emptyset$ and at every point $p\in E$:
\begin{enumerate}
\item[(GP1)] the active normals $\nor(p),a_1,\dots,a_k$ are linearly
independent;
\item[(GP2)] all pairwise angles among $\nor(p),a_1,\dots,a_k$ lie in
$[\theta_0,\pi-\theta_0]$; equivalently, the Gram matrix of the unit active
normals has smallest eigenvalue $\ge g_0=g_0(\theta_0,k)>0$.
\end{enumerate}
\end{definition}

Under (GP1)--(GP2) the map
$x\mapsto(\langle x,a_1\rangle,\dots,\langle x,a_k\rangle)$ is a submersion of
$\partial K$ near $E$, so $E$ is a compact $C^2$ submanifold of $\partial K$
of dimension
\[
j\;:=\;n-1-k\;\ge\;0 .
\]
Throughout we keep $0\le j\le n-1$; for $j=0$ (i.e.\ $k=n-1$) the stratum $E$
is a finite set of conical points and the exponent is the conical $1/n$ of
\cite{companion3,companion5}, so the interesting range is $1\le j\le n-1$,
i.e.\ $1\le k\le n-2$.  The orthogonal model of \cite{companion11} is the
special case $K=B^n$, $a_i=e_i$, $c_i=0$, for which $E=\{x\in S^{n-1}:x_1=\dots
=x_k=0\}$, $\nor(p)=p$, all active normals orthogonal, and any
$\theta_0\le\pi/2$ admissible.

We assume, as in all the companion papers, that $G$ has nonempty interior and
is uniformly thick at scale $r_0=r_0(K,h_k,\theta_0)>0$: there is a ball
$B(c^-,r_0)\subseteq G$.  Such a ball exists for any admissible configuration
because $E$ is a transversal intersection, so $G$ contains a fixed cone-cap
neighborhood pointing into $\bigcap_i\{\langle x,a_i\rangle<c_i\}$; we record
this as a standing hypothesis and absorb $r_0$ into the constants.

\subsection{The conjugate frame}

Fix $p\in E$.  The truncating normals $a_1,\dots,a_k$ are linearly
independent and, by (GP1), independent from $\nor(p)$.  We will need to write
the $k$ constraints in coordinates that diagonalize them while respecting the
tangent geometry at $p$.

\begin{definition}[Conjugate frame]\label{def:conj}
Let $V:=\mathrm{span}\{a_1,\dots,a_k\}$ and let
$\Pi:\R^n\to T_p\partial K=\nor(p)^\perp$ be the orthogonal projection onto the
tangent space at $p$.  Because $\nor(p)\notin V$ by (GP1), $\Pi|_V:V\to\Pi V$
is injective, so $\bar a_i:=\Pi a_i$ are linearly independent tangent vectors.
The \emph{conjugate frame} $f_1,\dots,f_k$ is the basis of $\Pi V\subseteq
T_p\partial K$ dual to $(\bar a_1,\dots,\bar a_k)$:
\[
\langle f_i,a_j\rangle=\delta_{ij},\qquad \langle f_i,\nor(p)\rangle=0
\qquad(1\le i,j\le k).
\]
\end{definition}

The two displayed conditions determine $f_1,\dots,f_k$ uniquely: they force
$f_i\in T_p\partial K$, and on $T_p\partial K$ the functional
$v\mapsto\langle v,a_j\rangle=\langle v,\bar a_j\rangle$ has the matrix
$\bar G=(\langle\bar a_i,\bar a_j\rangle)$, so $f_i=\sum_m(\bar G^{-1})_{im}
\bar a_m$.  By (GP2) the Gram matrix of $a_i$ has eigenvalues in
$[g_0,k]$, and the tangential projection along $\nor(p)$ changes lengths by a
factor bounded in terms of $\theta_0$ only, so
\begin{equation}\label{eq:fbound}
g_1(\theta_0,k)\;\le\;|f_i|\;\le\;g_2(\theta_0,k),
\qquad
|\langle f_i,f_j\rangle|\le g_2^2 ,
\end{equation}
with $g_1>0$.  In the new linear coordinates
\[
s_i\;:=\;\langle x-p,a_i\rangle\qquad(1\le i\le k),
\]
the truncation constraints $\langle x,a_i\rangle\le c_i$ read exactly
$s_i\le0$, since $\langle p,a_i\rangle=c_i$ for $p\in E$.  Thus the conjugate
frame diagonalizes the truncation: the $k$ half-spaces become the standard
orthant $\{s_1\le0,\dots,s_k\le0\}$, and $f_i$ is the direction in which only
$s_i$ moves (to first order, tangentially).  This is the device that replaces
the relation $\langle w,\tilde u\rangle=0$ used in the orthogonal model: there
$a_i=e_i$ is self-dual, $f_i=e_i$, and the windows are axis-parallel; in
general position the windows are spanned by the conjugate frame instead.

The invariant content we use repeatedly is the splitting of $T_p\partial K$
into the $k$-dimensional conjugate span $\Pi V$ and its $j$-dimensional
orthogonal complement
\[
T_pE\;=\;\{v\in T_p\partial K:\ \langle v,a_i\rangle=0\ \ \forall i\},
\qquad \dim T_pE=m:=j=n-1-k,
\]
the tangent space to the stratum; concrete orthonormal coordinates on each
factor are introduced as needed in \S\ref{ss:floor}.

\subsection{Main results}

The first result isolates the exponent.  For $p\in E$ let
\[
N(G,p)\;:=\;\cone\{\nor(p),a_1,\dots,a_k\}
\]
be the (closed, convex, simplicial, full in its $(k+1)$-dimensional span)
normal cone of $G$ at $p$, and for a unit vector $u$ in the relative interior
of $N(G,p)$ define the directional cap
\[
C_p^u(h)\;:=\;\{x\in G:\ \langle x,u\rangle\ \ge\ \langle p,u\rangle-h\}.
\]

\begin{proposition}[Universal cap exponent]\label{prop:exp}
Let $(K,h_k)$ be admissible and $p\in E$.  For every unit $u$ in the relative
interior of $N(G,p)$ there are $0<\gamma_1\le\gamma_2$ and $h_*>0$, depending
on $n,k,\theta_0$, the curvature band, and the angular distance of $u$ to
$\partial N(G,p)$, such that
\[
\gamma_1\,h^{\frac{n+1+k}2}\ \le\ \vol_n\!\bigl(C_p^u(h)\bigr)\ \le\
\gamma_2\,h^{\frac{n+1+k}2}
\qquad(0<h\le h_*),
\qquad \frac{n+1+k}2=\frac{2n-j}2 .
\]
The exponent is independent of the angles $\theta_0$ and of $u$ within the
cone.
\end{proposition}

\begin{theorem}[Section lemma, general position]\label{thm:A}
Let $(K,h_k)$ be admissible.  There are $a>0$ and $h_0>0$, depending only on
$n,k,\theta_0$ and the curvature band, such that for every unit vector $u$ and
every $0<h\le h_0$, writing $\tau_{\max}:=\max_{x\in G}\langle x,u\rangle$ and
$A_u(h):=\vol_{n-1}(G\cap\{\langle x,u\rangle=\tau_{\max}-h\})$,
\[
A_u(h)\ \ge\ a\,h^{\frac{n-1+k}2}.
\]
Consequently the UDC for the truncation set $h_k$ holds, with probability at
least $1-\frac1N$ for $N\ge N_0(n,k,\theta_0,K)$, at scaling factor
\[
\ell_N\;=\;C\,\Bigl(\frac{\log N}{N}\Bigr)^{2/(n+1+k)} .
\]
\end{theorem}

\begin{theorem}[Lower bound, general position]\label{thm:B}
Let $(K,h_k)$ be admissible.  There is $c>0$ such that with probability
$1-o(1)$ the doubling inequality fails for some slab truncated by $h_k$ at
every grid scale below $c\,(\log N/N)^{2/(n+1+k)}$; hence the UDC requires
$\ell\ge c\,(\log N/N)^{2/(n+1+k)}$.  The only configuration-dependent input
is the two-sided cap bound
\[
c_1\,h^{\frac{n+1+k}2}\ \le\ \mu\bigl(C_p(h)\bigr)\ \le\
c_2\,h^{\frac{n+1+k}2},
\qquad c_1=\omega(N_p)\cdot\kappa_{\mathrm{geom}}>0,
\]
where $\omega(N_p)$ is the normalized solid angle of the normal cone
$N(G,p)$ and $\kappa_{\mathrm{geom}}>0$ collects the curvature and
conjugate-frame constants.
\end{theorem}

\begin{corollary}[The interpolation law in general position]\label{cor:C}
Let $K$ be a $C^2$ strictly convex body with curvatures in
$[\lambda_{\min},\lambda_{\max}]$ and let $h_k$ be a general-position
simplicial truncation with dihedral angles bounded below by $\theta_0$, whose
lowest stratum $E$ has dimension $j=n-1-k$.  Then the critical scaling of the
uniform doubling condition is
\[
\ell_N\;=\;\Theta\!\Bigl(\bigl(\log N/N\bigr)^{2/(2n-j)}\Bigr),
\qquad 1\le j\le n-1,
\]
with exponent and logarithm exact and all constants depending only on
$n,k,\theta_0$ and the curvature band.  For $j=0$ (simplicial cones, $k=n-1$)
the conical scaling $\Theta((\log N/N)^{1/n})$ of \cite{companion3,companion5}
applies.  In particular the law holds for every simplicial cone and for every
nonsingular affine image of the orthogonal model.
\end{corollary}

Corollary~\ref{cor:C} is the general-position completion of
\cite[Theorem 1.1]{companion11}: it removes the roundness of the ball and the
orthogonality of the cuts, the only two special features of the model, while
keeping both halves of the bound and the constant in the logarithm.  The route
that does \emph{not} work --- straightening the configuration to the model by a
linear change of variables --- is discussed in \S\ref{sec:john}; the obstruction
there is what forces the intrinsic conjugate-frame argument of
\S\S\ref{sec:caps}--\ref{sec:upper}.

\paragraph{Organization.}
Section~\ref{sec:caps} sets up depth coordinates at the stratum and proves the
universal exponent (Proposition~\ref{prop:exp}) and the two-sided cap bound
(Lemma~\ref{lem:capmass}).  Section~\ref{sec:lower} proves
Theorem~\ref{thm:B}.  Section~\ref{sec:section} proves the section lemma
(Theorem~\ref{thm:A}) in general position, by induction on $k$ through the
slicing closure of the admissible class.  Section~\ref{sec:upper} assembles the
slab bound and the upper half of Theorem~\ref{thm:A}.  Section~\ref{sec:john}
records the John-position obstruction and the residual open case.

\section{Depth coordinates and the cap exponent}\label{sec:caps}

\subsection{Floor coordinates at a stratum point}\label{ss:floor}

Fix $p\in E$ and a unit vector $u$ in the relative interior of $N(G,p)$.  We
build coordinates adapted to $u$ in which the convex floor of $G$ near $p$ is
explicit.

Let $T:=T_p\partial K=\nor(p)^\perp$, split as $T=\Pi V\oplus T_pE$ with
$\Pi V=\mathrm{span}\{\bar a_1,\dots,\bar a_k\}$ (dimension $k$) and
$T_pE=\{v\in T:\langle v,a_i\rangle=0\}$ (dimension $m:=j=n-1-k$).  Write
$u=u^\nu\,\nor(p)+u^V+u^E$ with $u^\nu=\langle u,\nor(p)\rangle$,
$u^V\in\Pi V$, $u^E\in T_pE$.  Because $u\in\mathrm{relint}\,N(G,p)$ and
$N(G,p)=\cone\{\nor(p),a_1,\dots,a_k\}$, the expansion
$u=\beta_0\nor(p)+\sum_i\beta_ia_i$ has all $\beta_i>0$; since $u^E$ is the
component of $u$ orthogonal to $\nor(p)$ and to every $a_i$, and the $a_i$ span
$V\supseteq\Pi V$, one checks $u^E=0$.  Thus $u$ lies in the
$(k+1)$-plane $\mathrm{span}\{\nor(p)\}\oplus\Pi V$, the span of the cone, as it
must.

Choose orthonormal coordinates on $\R^n$ as follows.  The first coordinate is
$z:=\langle p-x,\nor(p)\rangle\ge0$, the depth below the tangent plane at $p$
(nonnegative on $K$ near $p$ by convexity).  On $T_pE$ choose an orthonormal
basis $\eta_1,\dots,\eta_m$ and set $y_a:=\langle x-p,\eta_a\rangle$.  On
$\Pi V$ we use the conjugate coordinates
$s_i:=\langle x-p,a_i\rangle$ of Definition~\ref{def:conj}; by
\eqref{eq:fbound} the change of basis between $(s_1,\dots,s_k)$ and an
orthonormal basis of $\Pi V$ has condition number bounded by
$g_2/g_1=:\Gamma(\theta_0,k)$, and we will only ever need this two-sided
comparison.  Thus $(z,y,s)\in\R\times\R^m\times\R^k$ is a system of
\emph{bounded-distortion} coordinates centered at $p$, with Jacobian bounded
above and below by constants depending on $\theta_0,k$.

\begin{lemma}[Convex floor]\label{lem:floor}
In the coordinates $(z,y,s)$ there are $r,L>0$ depending on $n,k,\theta_0$ and
the curvature band such that for $|y|,|s|\le r$, the boundary $\partial K$ is
the graph
\[
z\;=\;\Phi_0(y,s)\;:=\;\tfrac12\,Q(y,s)\;+\;O\!\bigl((|y|+|s|)^3\bigr),
\qquad
\lambda_{\min}\,(|y|^2+|s|^2)\ \le\ Q(y,s)\ \le\ \lambda_{\max}\,(|y|^2+|s|^2),
\]
where $Q$ is the second fundamental form $\mathrm{I\!I}_p$ read in these
coordinates, and the $O$-term has $C^0$ norm $\le L(|y|+|s|)^3$.  Moreover
$G\cap\{|y|,|s|\le r\}=\{z\ge\Phi_0(y,s)\}\cap\{s_1\le0,\dots,s_k\le0\}$ within
$\{|y|,|s|\le r\}$.
\end{lemma}

\begin{proof}
Since $\nor(p)$ is the outer normal, $\partial K$ is locally the graph of a
$C^2$ function $z=\phi(y,s)$ over $T$ with $\phi(0)=0$, $\nabla\phi(0)=0$, and
Hessian $\mathrm{I\!I}_p$, whose eigenvalues are the principal curvatures in
$[\lambda_{\min},\lambda_{\max}]$.  Taylor expansion gives $\Phi_0$ with the
stated quadratic bounds and a $C^2$ remainder, uniformly because the curvatures
are bounded and $K$ is $C^2$ with a modulus controlled on the compact stratum.
The constraint description holds because $\langle x-p,a_i\rangle=s_i$ exactly
and the curved faces $\{\langle x,a_i\rangle=c_i\}$ do not enter the
\emph{interior} normal stratum within the small box (they bound $G$ only where
some $s_i=0$, already included).
\end{proof}

We separate the directions of $u$ by whether the floor is flat or curved along
them.  By the previous paragraph $u=u^\nu\nor(p)+u^V$ with $u^\nu>0$ (else $u$
is tangent and not interior to the cone).  Write the slab depth
$\langle p-x,u\rangle$, i.e.\ the cap function, in coordinates.  Using
$x-p=-z\,\nor(p)+\sum_a y_a\eta_a+\Pi(x-p)$ and that $u^E=0$,
\begin{align}
\langle p-x,u\rangle
&=u^\nu z-\langle x-p,u^V\rangle \notag\\
&=u^\nu z-\sum_{i=1}^k\beta_i\,\langle x-p,a_i\rangle
\;=\;u^\nu z-\sum_{i=1}^k\beta_i\,s_i,
\label{eq:capfun}
\end{align}
because $u^V=\sum_i\beta_i\bar a_i$ pairs with $x-p$ through
$\langle\cdot,a_i\rangle=s_i$ on $T$ (the $\nor(p)$-component of $x-p$ is
orthogonal to $\Pi V$).  Here $\beta_i>0$ are the cone coefficients of $u$, and
by (GP2) and \eqref{eq:fbound} they obey $\beta_i\in[\beta_-,\beta_+]$ with
$\beta_->0$ depending on the angular distance of $u$ to $\partial N(G,p)$.

\subsection{Proof of Proposition~\ref{prop:exp}}

\begin{proof}
On $C_p^u(h)$ we have $\langle p-x,u\rangle\le h$, $s_i\le0$, and
$z\ge\Phi_0(y,s)$.  Substitute \eqref{eq:capfun}: with $\sigma_i:=-s_i\ge0$,
\begin{equation}\label{eq:capset}
C_p^u(h)\cap\{|y|,|s|\le r\}
=\Bigl\{(y,\sigma,z):\ \sigma_i\ge0,\ z\ge\Phi_0(y,-\sigma),\
u^\nu z+\sum_i\beta_i\sigma_i\le h\Bigr\}.
\end{equation}
The cap is contained in this box for $h\le h_*$ small, because the cap radius
is $O(\sqrt h)$ (shown below) and then $\Phi_0\ge\lambda_{\min}(|y|^2+|s|^2)/2$
forces $|y|,|s|\le r$.

\emph{Rescaling.}  Put $y=\sqrt h\,Y$, $\sigma=h\,S$ (note: linear in $h$ for
the truncated directions, square-root for the tangent-to-stratum directions),
and $z=h\,Z$.  Then $dy\,d\sigma\,dz=h^{m/2}\cdot h^{k}\cdot h\;dY\,dS\,dZ$, and
\[
\Phi_0(y,-\sigma)=\tfrac12 Q(\sqrt h Y,-hS)+O(h^{3/2})
=\tfrac{h}{2}Q(Y,0)+O(h^{3/2}),
\]
because $Q(\sqrt h Y,-hS)=h\,Q(Y,0)-h^{3/2}(\text{cross})+h^2Q(0,S)$ and the
last two terms are $o(h)$.  Hence after dividing the cap inequality by $h$,
\begin{equation}\label{eq:rescaled}
C_p^u(h)\ \longleftrightarrow\
\Bigl\{(Y,S,Z):\ S_i\ge0,\
Z\ \ge\ \tfrac12 Q(Y,0)+o(1),\
u^\nu Z+\sum_i\beta_i S_i\ \le\ 1+o(1)\Bigr\}.
\end{equation}
The limiting region (as $h\to0$) is
\[
R:=\Bigl\{(Y,S,Z):\ S\ge0,\ Z\ge\tfrac12 Q(Y,0),\
u^\nu Z+\textstyle\sum_i\beta_iS_i\le1\Bigr\},
\]
a fixed set of finite, positive $n$-volume:
finiteness because $Q(Y,0)\ge\lambda_{\min}|Y|^2$ confines $Y$ to
$|Y|\le\sqrt{2/(u^\nu\lambda_{\min})}$ and then $Z\le1/u^\nu$, $S_i\le
1/\beta_i$; positivity because the relative interior is nonempty (take $Y=0$,
small $S_i>0$, $Z$ slightly above $0$).  The $o(1)$ perturbations are uniform
on compact sets and the regions are nested up to $1\pm o(1)$ dilations, so by
dominated convergence
\[
\vol_n\!\bigl(C_p^u(h)\bigr)
=h^{\,1+k+m/2}\bigl(\vol_n(R)+o(1)\bigr),
\qquad 1+k+\tfrac m2=\frac{(n-1)+1+k}2=\frac{n+1+k}2 .
\]
The exponent counts $1$ for the depth $z$ (length $h$), $1$ for each truncated
direction $s_i$ (length $h$, a \emph{linear} floor $\beta_i\sigma_i$), and
$\tfrac12$ for each of the $m=j$ stratum directions $y_a$ (length $\sqrt h$, a
\emph{quadratic} floor $\tfrac12 Q$).  None of these scalings sees the angles,
only the dichotomy linear-vs-quadratic floor; this proves the
angle-independence and gives $\gamma_1,\gamma_2$ as $\vol_n(R)(1\mp o(1))$.
Finally the cap-radius claim: on $R$, $|Y|,|S|,|Z|=O(1)$, so on $C_p^u(h)$,
$|y|=O(\sqrt h)$, $|s|=O(h)$, $z=O(h)$, hence $|x-p|=O(\sqrt h)$, justifying the
box confinement for $h\le h_*$.
\end{proof}

\begin{remark}\label{rem:why}
Proposition~\ref{prop:exp} is the conceptual core.  In the orthogonal model
the same computation has $a_i=e_i$, $\beta_i$ the positive coordinates of $u$,
and $Q=\mathrm{Id}$ (the round sphere), so the floor is
$\tfrac12|Y|^2+\sum\beta_i\sigma_i$ and one reads off $(n+1+k)/2$ directly.
General position changes $Q$ to an anisotropic but uniformly elliptic form and
changes the truncated directions from coordinate axes $e_i$ to the conjugate
directions $f_i$; neither change touches the count of linear versus quadratic
floor directions, which is all the exponent records.  This is why no rounding
of $K$ or orthogonalization of the $a_i$ is needed for the exponent --- and,
as \S\ref{sec:john} shows, no such rounding exists.
\end{remark}

\subsection{The stratum cap and its mass}\label{ss:capmass}

For the lower bound we need a canonical cap at each $p\in E$ and its mass with
the correct constant.  Let
\[
u_p\;:=\;\frac{\nor(p)+\sum_{i=1}^k a_i}{\bigl|\nor(p)+\sum_i a_i\bigr|},
\qquad
C_p(h)\;:=\;C_p^{u_p}(h)
=\{x\in G:\ \langle x,u_p\rangle\ge\langle p,u_p\rangle-h\}.
\]
By (GP1)--(GP2) the vector $\nor(p)+\sum_i a_i$ is a fixed positive combination
of the cone generators, hence $u_p\in\mathrm{relint}\,N(G,p)$ with angular
distance to $\partial N(G,p)$ bounded below by a function of $\theta_0,k$;
$\langle p,u_p\rangle=\max_{x\in G}\langle x,u_p\rangle$ because $u_p\in
N(G,p)$ (supporting at $p$).  So $C_p(h)$ is a cap of the type in
Proposition~\ref{prop:exp}, uniformly over $E$.

\begin{lemma}[Stratum cap mass and footprint]\label{lem:capmass}
There are $0<c_1\le c_2$ and $h_1>0$ depending on $n,k,\theta_0$ and the
curvature band, and $c_3=c_3(\theta_0,k)$, such that for all $p\in E$ and
$0<h\le h_1$,
\[
c_1\,h^{\frac{n+1+k}2}\ \le\ \mu\bigl(C_p(h)\bigr)\ \le\
c_2\,h^{\frac{n+1+k}2},
\qquad
C_p(h)\subseteq\{x:\ |x-p|\le c_3\sqrt h\}.
\]
Moreover one may take $c_1=\omega(N_p)\cdot\kappa_{\mathrm{geom}}$, where
$\omega(N_p)\in(0,1]$ is the normalized solid angle of $N(G,p)$ and
$\kappa_{\mathrm{geom}}>0$ depends only on the curvature band and
$\theta_0,k$.
\end{lemma}

\begin{proof}
Apply Proposition~\ref{prop:exp} with $u=u_p$: it gives
$\vol_n(C_p(h))=h^{(n+1+k)/2}(\vol_n(R_p)+o(1))$ and the footprint
$|x-p|\le c_3\sqrt h$, and dividing by $\vol_n(K)$ gives the measure bound with
$c_2=2\vol_n(R_p)/\vol_n(K)$, $c_1=\tfrac12\vol_n(R_p)/\vol_n(K)$ for
$h\le h_1$, the constants uniform in $p$ because $R_p$ depends continuously on
the compact data $(\nor(p),a_i,\mathrm{I\!I}_p)$ and is bounded away from
degenerate.

For the solid-angle form of $c_1$, compute $\vol_n(R_p)$ from below by the
cone factorization.  In \eqref{eq:capset}, drop the curvature in $Z$
(i.e.\ restrict to $Z\ge0$ but allow $Z\ge\tfrac12Q(Y,0)$); the region $R_p$
contains the product slab
\[
\{Y:\tfrac12 Q(Y,0)\le\tfrac1{2u^\nu}\}\times
\{(S,Z):\ S\ge0,\ Z\ge0,\ u^\nu Z+\textstyle\sum_i\beta_i S_i\le\tfrac12\},
\]
whose $Y$-factor has volume
$\ge\det(\mathrm{I\!I}_p)^{-1/2}\,\vol_m(\{|\cdot|\le\rho\})\ge
\lambda_{\max}^{-m/2}\,\omega_m\rho^m$ (an ellipsoid with semi-axes
$\ge(u^\nu\lambda_{\max})^{-1/2}$), and whose $(S,Z)$-factor is the simplex
$\{S\ge0,Z\ge0,\ u^\nu Z+\sum\beta_iS_i\le\tfrac12\}$ of volume
$\tfrac1{(k+1)!}\,(\tfrac12)^{k+1}/(u^\nu\prod_i\beta_i)$.  Now
$u^\nu\prod_i\beta_i$ is, up to the bounded Gram determinant of the cone
generators, exactly the reciprocal of the normalized solid angle of
$N(G,p)=\cone\{\nor(p),a_i\}$ at unit radius: indeed the simplex
$\{\beta_0\ge0,\beta_i\ge0,\ \beta_0u^\nu+\sum\beta_i\beta_i\le\tfrac12\}$ in
cone-coordinates has Euclidean volume proportional to
$\sqrt{\det \mathrm{Gram}(\nor(p),a_i)}$, and that determinant is bounded in
$[g_0^{(k+1)/2}, (k+1)^{(k+1)/2}]$ by (GP2).  Collecting the bounded curvature
and Gram factors into $\kappa_{\mathrm{geom}}$ yields
$\vol_n(R_p)\ge\omega(N_p)\,\kappa_{\mathrm{geom}}\,\vol_n(K)$, hence
$c_1=\omega(N_p)\kappa_{\mathrm{geom}}>0$.  Positivity of $\omega(N_p)$ is
exactly (GP1): a degenerate (lower-dimensional, $\omega=0$) cone would mean
$\nor(p),a_i$ linearly dependent.
\end{proof}

\begin{remark}\label{rem:reflfrac}
This is where the orthogonal model's exact fraction $2^{-k}$ is replaced.
There, $N(G,p)=\cone\{p,e_1,\dots,e_k\}$ is a right-angled corner, its
normalized solid angle is exactly $2^{-k}$ of the half-space solid angle, and
the all-negative orthant keeps that fraction by reflection symmetry
(\cite[Lemma 2.2]{companion11}).  In general position the corner is oblique,
the reflection symmetry is gone, and the surviving fraction of the cap is the
genuine normalized solid angle $\omega(N_p)$ of the oblique cone --- still a
fixed positive number by (GP2), but no longer a clean power of two.  Only
positivity, not the exact value, is needed downstream.
\end{remark}

\begin{lemma}[Stall]\label{lem:stall}
Fix $\kappa>1$, $\sigma>0$, $\delta\ge\frac3N$, $p\in E$ and
$0<h\le\frac H4$, $H\le h_1$.  If exactly one sample point lies in $C_p(H)$ and
it lies in $C_p(h)$, then for every $w\le H/\kappa$ the slab with direction
$u_p$, center that sample point, and truncation set $h_k$ violates the doubling
inequality; hence the UDC fails for every $\ell\le H/\kappa^2$.
\end{lemma}

\begin{proof}
Verbatim the companion stall lemma \cite[Lemma 2.4]{companion11}: for
$w'\le H$ the slab is contained in $C_p(h+\frac{w'}2)\subseteq C_p(H)$, so the
empirical masses of $S_w$ and $S_{\kappa w}$ both equal $\frac1N<
\min\{\delta,(1+\sigma)\frac1N\}$.  Nothing in the argument uses orthogonality
or roundness; it uses only that caps are nested and that $C_p(H)$ holds a
single point.
\end{proof}

\section{The lower bound}\label{sec:lower}

\begin{proof}[Proof of Theorem~\ref{thm:B}]
This is the second-moment scheme of \cite{companion7,companion11} verbatim,
fed by Lemma~\ref{lem:capmass}.  Set
$q:=(4\kappa^2)^{(n+1+k)/2}\,c_2/c_1$,
$\varepsilon_3:=\frac{j}{4q(2n-j)}$, $\tau_N:=\varepsilon_3\log N$.  Define
$h_N$ by $\mu(C_p(h_N))=\tau_N/N$; this is well defined and, by
Lemma~\ref{lem:capmass}, $h_N\asymp(\tau_N/N)^{2/(n+1+k)}$ uniformly in $p$
(the constants $c_1,c_2$ are $p$-independent).  Put $H_N:=4\kappa^2h_N$, so
$\mu(C_p(H_N))\le q\,\tau_N/N$.

By the footprint bound the caps have radius $\le c_3\sqrt{H_N}$ around their
centers $p\in E$.  Since $E$ is a compact $C^2$ $j$-manifold, a maximal
$2c_3\sqrt{H_N}$-separated family $p_1,\dots,p_M\in E$ has
\[
M\ \ge\ c\,\bigl(\sqrt{H_N}\bigr)^{-j}\ \ge\
c'\,(N/\tau_N)^{\frac{j}{n+1+k}},
\qquad \frac{j}{n+1+k}=\frac{j}{2n-j},
\]
because the $j$-dimensional volume of $E$ is a fixed positive constant and a
metric ball of radius $\rho$ in $E$ has volume $\asymp\rho^j$ (curvatures of
$E$ are bounded, by admissibility).  The outer caps $C_{p_\alpha}(H_N)$ are
pairwise disjoint: if $x\in C_{p_\alpha}(H_N)\cap C_{p_\beta}(H_N)$ then
$|x-p_\alpha|,|x-p_\beta|\le c_3\sqrt{H_N}$, so
$|p_\alpha-p_\beta|\le2c_3\sqrt{H_N}$, contradicting separation.

Let $E_\alpha$ be the event that $C_{p_\alpha}(H_N)$ contains exactly one
sample point and that point lies in $C_{p_\alpha}(h_N)$, and
$X:=\sum_\alpha\mathbf1_{E_\alpha}$.  Exactly as in the companion,
\[
\PP(E_\alpha)=N\,\mu(C_{p_\alpha}(h_N))\,(1-\mu(C_{p_\alpha}(H_N)))^{N-1}
\ge\tau_N\,N^{-2q\varepsilon_3},
\]
\[
\E X=\sum_\alpha\PP(E_\alpha)\ge c'\,N^{\frac j{2n-j}-2q\varepsilon_3}\,
\tau_N^{\,1-\frac j{2n-j}}\ \ge\ c'\,N^{\frac{j}{2(2n-j)}}\to\infty,
\]
using $2q\varepsilon_3=\frac{j}{2(2n-j)}$.  The outer caps being disjoint, the
events $E_\alpha$ are negatively associated up to the standard correction, and
the variance ratio is $\exp(O(N\,\mu(C(H_N))^2))=1+o(1)$; Chebyshev gives
$\PP(X=0)\le(\E X)^{-1}+o(1)=o(1)$.  On $\{X\ge1\}$, Lemma~\ref{lem:stall}
forces $\ell>H_N/\kappa^2=4h_N\ge c\,(\log N/N)^{2/(n+1+k)}$.
\end{proof}

\section{The section lemma in general position}\label{sec:section}

We prove Theorem~\ref{thm:A}.  As in \cite{companion11} it is equivalent, and
more convenient, to prove the cap-mass lower bound
\begin{equation}\label{eq:capgoal}
\Lambda_u(h):=\mu\bigl(\{x\in G:\ \langle x,u\rangle\ge\tau_{\max}-h\}\bigr)
\ \ge\ b\,h^{\frac{n+1+k}2}
\qquad(0<h\le h_0)
\end{equation}
for every unit $u$, together with a bulk anchor; the conversion to sections is
identical to the model and is recalled at the end.  The proof is by induction
on the number $k$ of truncations.  The induction is legitimate because the
admissible class is closed under the slicing used below
(Lemma~\ref{lem:slice}).

\subsection{Closure under slicing}

\begin{lemma}[Slices of admissible configurations are admissible]\label{lem:slice}
Let $(K,h_k)$ be admissible with $k\ge1$, fix $i_0\in\{1,\dots,k\}$ and a
level $s\in\R$ near $c_{i_0}$, and set
\[
K_s\;:=\;K\cap\{\langle x,a_{i_0}\rangle=s\},
\qquad
h_{k-1}^{(s)}\;:=\;\{(a_i,c_i):i\ne i_0\}\big|_{K_s}.
\]
For $s$ in a one-sided neighborhood of $c_{i_0}$ on which $K_s\ne\emptyset$,
the pair $(K_s,h_{k-1}^{(s)})$ is an admissible $(n-1,k-1)$ configuration with
the same angular constant $\theta_0'\ge\theta_0/2$ and curvature band
$[\lambda_{\min},\lambda_{\max}']$, $\lambda_{\max}'\le\lambda_{\max}/\sin
\theta_0$, and its stratum $E_s=E\cap\{\langle x,a_{i_0}\rangle=s\}$ has
dimension $j$ (one less ambient dimension, one fewer constraint).
\end{lemma}

\begin{proof}
$K_s$ is the intersection of the $C^2$ strictly convex body $K$ with an affine
hyperplane $\{\langle\cdot,a_{i_0}\rangle=s\}$ transversal to $\partial K$ near
$E$ (transversality is (GP1): $a_{i_0}$ is independent of $\nor$ along $E$, so
the hyperplane is not tangent).  A transversal hyperplane section of a $C^2$
strictly convex body is a $C^2$ strictly convex body in the hyperplane, with
principal curvatures bounded below by $\lambda_{\min}$ and above by
$\lambda_{\max}/\sin\theta(a_{i_0},\nor)\le\lambda_{\max}/\sin\theta_0$
(the section curvature is the normal curvature divided by the sine of the
incidence angle).  The remaining normals $a_i$ ($i\ne i_0$), projected into the
hyperplane $a_{i_0}^\perp$, are the truncating normals of the slice; their
mutual angles and their angles to the section normal stay in $[\theta_0/2,\pi
-\theta_0/2]$ because projecting along $a_{i_0}$ (which makes angle $\ge\theta_0$
with each of them and with $\nor$) distorts angles by at most $\theta_0/2$ for
$s$ close enough to $c_{i_0}$.  Linear independence (GP1) of
$\{\nu_{K_s},a_i:i\ne i_0\}$ in the hyperplane follows from (GP1) for the full
list minus $a_{i_0}$.  The stratum identity is immediate.
\end{proof}

\begin{remark}
Lemma~\ref{lem:slice} is the general-position analogue of the fact, used freely
in \cite{companion11}, that a coordinate slice of the orthogonal model
$B^n\cap\{x_i\le0\}$ is the orthogonal model $B^{n-1}\cap\{x_i\le0\}$ in one
lower dimension.  There the slice is again a round ball cut by coordinate
half-spaces; here it is an oblique section, a $C^2$ convex body cut by oblique
half-spaces, but still admissible.  This closure is what lets the induction run
without ever leaving the class.
\end{remark}

\subsection{Proof of the cap-mass bound \eqref{eq:capgoal}}

\begin{proof}
We induct on $k$.  The base case $k=0$ is the smooth strictly convex body: by
Lemma~\ref{lem:floor} with no constraints, the cap of depth $h$ in direction
$u=\nor(p_u)$ (with $p_u$ the unique boundary maximizer of $\langle\cdot,u
\rangle$) has, after the rescaling $y=\sqrt h Y$, $z=hZ$,
$\mu(C(h))\ge c(n,\lambda_{\max})h^{(n+1)/2}$, the standard cap bound for a
body with bounded curvature \cite{companion3}.  Assume \eqref{eq:capgoal} for
all admissible $(n-1,k-1)$ configurations.

Fix an admissible $(K,h_k)$ and a unit $u$.  Let
$p_u\in\partial K$ maximize $\langle\cdot,u\rangle$ over $G$; by convex duality
$p_u$ lies in the stratum whose normal cone contains $u$.  Let $S\subseteq
\{1,\dots,k\}$ be the active set at $p_u$, i.e.\ the constraints with
$\langle p_u,a_i\rangle=c_i$.  We distinguish three cases by the position of
$u$ relative to the normal cone $N(G,p_u)=\cone\{\nor(p_u),a_i:i\in S\}$.

\medskip
\emph{Case 1 ($u$ interior to the full normal cone; $S=\{1,\dots,k\}$ and $u\in
\mathrm{relint}\,N(G,p_u)$).}
Then $p_u\in E$ and Proposition~\ref{prop:exp} (via Lemma~\ref{lem:capmass}'s
proof, which only used interiority and admissibility) gives directly
\[
\Lambda_u(h)=\mu(C_{p_u}^u(h))\ \ge\ \gamma_1\,h^{(n+1+k)/2},
\]
with $\gamma_1>0$ depending on the angular distance of $u$ to
$\partial N(G,p_u)$, hence on $\theta_0,k$ uniformly.  This is
\eqref{eq:capgoal}.

\medskip
\emph{Case 2 (some constraints inactive at $p_u$; $S\subsetneq\{1,\dots,k\}$).}
Then near $p_u$ the constraints $\{i\notin S\}$ are slack:
$\langle p_u,a_i\rangle<c_i$ for $i\notin S$, so there is $\rho>0$ with
$\langle x,a_i\rangle<c_i$ for all $i\notin S$ and all $x$ with
$|x-p_u|\le\rho$.  For $h\le h_0(\rho)$ the cap $C^u(h)$ has radius
$\le c_3\sqrt h\le\rho$ (footprint bound), hence lies in
$\{x:\langle x,a_i\rangle<c_i,\ i\notin S\}$, where $G$ agrees with the
sub-configuration $G_S:=K\cap\bigcap_{i\in S}\{\langle x,a_i\rangle\le c_i\}$.
Now $(K,h_{|S|})$ with truncation set $\{(a_i,c_i):i\in S\}$ is an admissible
$(n,|S|)$ configuration with $|S|\le k$ active constraints and the \emph{same}
$\theta_0$ (a sub-list of admissible normals is admissible), and $u\in
\mathrm{relint}\,N(G_S,p_u)$ with $p_u$ on \emph{its} lowest stratum
$E_S\supseteq E$.  If $|S|<k$ apply the induction hypothesis in the form
already proved for fewer constraints (Case 1 for $G_S$ when $u$ is interior to
$N(G_S,p_u)$, else recurse into Case 2/3 for $G_S$); the exponent it yields is
$(n+1+|S|)/2\le(n+1+k)/2$, and since $h\le1$,
\[
\Lambda_u(h)=\mu(C^u(h)\cap G_S)\ \ge\ b_S\,h^{(n+1+|S|)/2}\ \ge\
b_S\,h^{(n+1+k)/2}.
\]
The bound $h\le1$ only helps; we lose nothing because the larger exponent is
the worst case and it is achieved exactly when all $k$ constraints are active,
i.e.\ in Cases 1 and 3.

\medskip
\emph{Case 3 ($u$ near-tangent to an active facet; $S=\{1,\dots,k\}$ but $u\in
\partial N(G,p_u)$ or just inside, near the wall $\{\beta_{i_0}\to0\}$).}
This is the inductive heart.  Suppose $u$ is within angle $\theta_0/4$ of the
facet of $N(G,p_u)$ opposite the generator $a_{i_0}$, i.e.\ the cone
coefficient $\beta_{i_0}$ of $u$ is the smallest and may be near $0$.  We slice
out $a_{i_0}$.  For $s\in[\,c_{i_0}-\tfrac{h}{2|u^{(i_0)}|},\,c_{i_0}]$ ---
an interval of length $\tfrac{h}{2|u^{(i_0)}|}\ge\tfrac h2$, where
$u^{(i_0)}:=\langle u,a_{i_0}\rangle$ is bounded below by a constant times
$\beta_{i_0}$ but the interval length uses only $|u^{(i_0)}|\le1$ --- consider
the slice
\[
K_s=K\cap\{\langle x,a_{i_0}\rangle=s\},
\qquad
G_s=G\cap\{\langle x,a_{i_0}\rangle=s\}
=K_s\cap\bigcap_{i\ne i_0}\{\langle x,a_i\rangle\le c_i\}.
\]
By Lemma~\ref{lem:slice}, $(K_s,h_{k-1}^{(s)})$ is admissible of type
$(n-1,k-1)$.  On the slice, write
$\langle x,u\rangle=u^{(i_0)}s+\langle x,u_\perp\rangle$, where
$u_\perp$ is the orthogonal projection of $u$ onto $a_{i_0}^\perp$,
$|u_\perp|\le1$.  The unit direction $u_\perp/|u_\perp|$ is admissible for the
slice (it is some unit vector in the hyperplane), and the maximizer of
$\langle\cdot,u_\perp\rangle$ over $G_s$ approaches $p_u$ as $s\to c_{i_0}$; the
maximum value $\tau_{\max}^{(s)}$ differs from $\tau_{\max}-u^{(i_0)}s$ by at
most $O(h)$ for $s$ within $O(h)$ of $c_{i_0}$, by the $C^2$ dependence of the
section on $s$ (curvatures bounded via Lemma~\ref{lem:slice}).  Hence the cap
of depth $h$ for $(G,u)$ contains, over each such $s$, the image of the cap of
depth $h_*:=h/(4|u_\perp|)\ge h/4$ for $(G_s,u_\perp/|u_\perp|)$, once we budget
the $\le2h$ of depth spent on the $s$-displacement and the section drift within
a total depth $4h$ (and then replace $h$ by $h/4$ throughout, which only changes
constants).  By the induction hypothesis on $(n-1,k-1)$ at depth
$\min(h_*,h_0)$ --- if $h_*>h_0$ use monotonicity $\Lambda(h_*)\ge\Lambda(h_0)
\ge b\,h_0^{(n+k-1)/2}\ge b'h^{(n+k-1)/2}$ --- the slice cap has
$(n-1)$-mass $\ge b(n-1,k-1)\,h^{((n-1)+1+(k-1))/2}=b\,h^{(n+k-1)/2}$.
Integrating over the $s$-interval of length $\ge h/2$ (the co-area factor
$|\nabla\langle\cdot,a_{i_0}\rangle|=1$ is constant) by Fubini,
\[
\Lambda_u(h)\ \ge\ \frac h2\cdot b(n-1,k-1)\,h^{\frac{n+k-1}2}
\;=\;b'\,h^{1+\frac{n+k-1}2}\;=\;b'\,h^{\frac{n+1+k}2}.
\]
This proves \eqref{eq:capgoal} in Case 3 and closes the induction.

\medskip
The three cases are exhaustive: either all $k$ constraints are active at $p_u$
(Cases 1 and 3, split by whether $u$ is interior or near-tangent to the cone)
or some are inactive (Case 2).  In every case
$\Lambda_u(h)\ge b\,h^{(n+1+k)/2}$ with $b=b(n,k,\theta_0,\lambda_{\min},
\lambda_{\max})>0$, proving \eqref{eq:capgoal}.
\end{proof}

\subsection{Conversion to sections}

Exactly as in \cite[\S4]{companion11}.  By Brunn--Minkowski the section area
$A_u(\cdot)$ is unimodal in the depth.  The thickness ball $B(c^-,r_0)\subseteq
G$ provides a bulk depth $h^\dagger\ge c(\theta_0,k)>0$ at which
$A_u(h^\dagger)\ge a_1>0$ (its $\langle\cdot,u\rangle$-range is an interval of
length $2r_0$ inside the range of $G$, with central sections of radius
$\ge\frac{\sqrt3}2 r_0$).  For $h\le h_0\le h^\dagger$: below the mode of
$A_u$, $A_u$ is nondecreasing on $[0,h]$ and $A_u(h)\ge\Lambda_u(h)/h\ge
b\,h^{(n-1+k)/2}$ by \eqref{eq:capgoal}; above the mode, unimodality gives
$A_u(h)\ge A_u(h^\dagger)\ge a_1\ge a_1h^{(n-1+k)/2}$.  Either way
$A_u(h)\ge a\,h^{(n-1+k)/2}$, which is Theorem~\ref{thm:A}.

\begin{remark}\label{rem:device}
The witness in Case 1 is the conjugate-frame box: take $Y$, $S$ in the limiting
region $R$ of Proposition~\ref{prop:exp} and pull back to a box in $x$ whose
$f_i$-edges have length $\asymp h$ (truncated directions) and whose
$\eta_a$-edges have length $\asymp\sqrt h$ (stratum directions), capped by the
fibre of length $\asymp h$ in the $u$-direction.  The orthogonal model's box
\cite[\S4]{companion11} is the special case $f_i=e_i$; the only structural
change is that the truncated edges point along the conjugate directions
$f_i\in\mathrm{span}\{a_i\}$ rather than the coordinate axes, which is forced
because the constraints are diagonalized by $s_i=\langle x-p,a_i\rangle$, not by
the standard coordinates.  Case 3 is the mechanism by which face strata inherit
the law from lower-dimensional admissible models, one truncation per dimension,
through the slicing closure of Lemma~\ref{lem:slice}.
\end{remark}

\section{Slab bound and the upper bound}\label{sec:upper}

\begin{theorem}[Slab volume]\label{thm:slab}
There are $c>0$ and $L_0>0$, depending only on $n,k,\theta_0$ and the curvature
band, with
\[
\vol_n\Bigl(G\cap\{|\langle y-x,u\rangle|\le L\}\Bigr)\ \ge\
c\,L^{\frac{n+1+k}2}
\qquad(x\in G,\ u\in S^{n-1},\ 0<L\le L_0).
\]
\end{theorem}

\begin{proof}
Anchor--unimodality as in \cite{companion3,companion11}.  Apply
Theorem~\ref{thm:A} to $u$ and to $-u$: there are anchors at depth $L/4$ from
each $u$-extreme of $G$ with section value $\ge a\,(L/4)^{(n-1+k)/2}$.  By
Brunn--Minkowski unimodality the section area exceeds the smaller anchor value
throughout the middle of the $\langle\cdot,u\rangle$-range, whose complement has
length $L/2$; the range has length $\ge2r_0\ge2L$ by thickness; and the slab
keeps length $\ge L$ inside the range, of which $\ge L/2$ is in the middle.
Integrating the section area over that middle length gives
$\ge\frac L2\cdot a(L/4)^{(n-1+k)/2}=c\,L^{(n+1+k)/2}$.
\end{proof}

\begin{proof}[Proof of Theorem~\ref{thm:A}, upper bound]
Insert Theorem~\ref{thm:slab} into the high-probability skeleton of
\cite{companion2,companion3}, identical to \cite[\S5]{companion11}.  The
population dichotomy holds with thickness radius $r_0$; the VC relative-deviation
bounds apply to slabs cut by the $k$ half-spaces of $h_k$ (a slab truncated by
$h_k$ is an intersection of $k+2$ half-spaces, VC dimension $O(nk)$, so the
uniform deviation $u_N\asymp\log N/N$ controls all of them simultaneously); and
a sample point in $S_w$ yields $\mu(S_{\kappa w})\ge32u_N$ for every
$w\ge\ell_N:=C\,(u_N)^{2/(n+1+k)}\asymp(\log N/N)^{2/(n+1+k)}$, by
Theorem~\ref{thm:slab} with $L=\kappa w$.  This proves the UDC at scaling
$\ell_N$ with probability $\ge1-\frac1N$.
\end{proof}

\noindent Theorems~\ref{thm:A} and \ref{thm:B} together give
Corollary~\ref{cor:C}.

\section{Why rounding fails, and the open case}\label{sec:john}

\subsection{The John-position obstruction}

The shortest conceivable proof of Corollary~\ref{cor:C} would map the admissible
configuration to the orthogonal model by a single affine change of variables
and quote \cite{companion11}.  This is provably impossible already at the
linear level, and the obstruction is instructive.

\begin{proposition}[No simultaneous rounding and orthogonalization]\label{prop:john}
Let $a_1,\dots,a_k\in S^{n-1}$ ($k\ge2$) be unit normals, not pairwise
orthogonal.  There is no invertible linear map $A\in GL_n(\R)$ such that
$A(B^n)$ is a ball and $Aa_1,\dots,Aa_k$ are pairwise orthogonal.
\end{proposition}

\begin{proof}
$A(B^n)$ is a ball iff $A^{\mathsf T}A=\lambda I$ for some $\lambda>0$ (then
$A=\sqrt\lambda\,O$, $O$ orthogonal).  But then for all $i,j$,
\[
\langle Aa_i,Aa_j\rangle=\langle A^{\mathsf T}Aa_i,a_j\rangle
=\lambda\,\langle a_i,a_j\rangle,
\]
so the inner products are scaled by the common factor $\lambda$, and the
\emph{angles} between the $a_i$ are preserved.  In particular $Aa_i\perp Aa_j$
forces $a_i\perp a_j$.  If some pair $a_i,a_j$ is non-orthogonal, no such $A$
orthogonalizes it while rounding the ball.
\end{proof}

Thus one cannot have both: any linear map that orthogonalizes the cuts
distorts $B^n$ into a genuine ellipsoid (and, for a general $K$, distorts the
curvature band in a direction-dependent way), and any linear map that keeps the
body a ball preserves all the cut angles.  Affinely, $K$ and the half-spaces are
\emph{jointly} rigid up to a similarity: the configuration carries a true
modulus, the angles, that no volume-preserving or shape-preserving map can
delete.  The argument of \S\S\ref{sec:caps}--\ref{sec:upper} is therefore
intrinsic of necessity: the conjugate frame diagonalizes the constraints
\emph{without} touching $K$, paying for the obliqueness with the anisotropic but
uniformly elliptic floor $Q$ of Lemma~\ref{lem:floor} and the solid angle
$\omega(N_p)$ of Lemma~\ref{lem:capmass}, rather than reducing to a right-angled
model.  This is also why the exponent, which depends only on the
linear-versus-quadratic floor dichotomy (Remark~\ref{rem:why}), survives
unchanged while the constants acquire $\theta_0$-dependence.

\begin{remark}
For $k=1$ a single cut can always be orthogonalized trivially, and indeed a
similarity reduces one oblique cut of a ball to the model; the obstruction is
genuinely a $k\ge2$ phenomenon, the first place where two angles must coexist.
This matches the fact that the orthogonal-model induction of
\cite{companion11} and of \S\ref{sec:section} only becomes nontrivial at
$k\ge2$.
\end{remark}

\subsection{Curved truncating surfaces and the companion}

Admissibility (Definition~\ref{def:adm}) keeps the truncating boundaries
\emph{flat} (half-spaces $\langle x,a_i\rangle\le c_i$).  The curved companion
\cite{companion15curved} treats the next case, where each flat constraint is
replaced by a \(C^2\) truncating surface \(g_i(x)\le0\), under the same
transversality philosophy.  For each fixed curved configuration it proves that
the exponent survives:
\[
\ell_N\;\asymp\;\Bigl(\frac{\log N}{N}\Bigr)^{2/(2n-j)}.
\]
The remaining issue is not the exponent but uniformity of the constants across
all curvature-bounded families.  With curved truncations the second fundamental
forms of the truncating walls enter the same rescaled floor block as the
stratum form, so the flat factorization
\(c_1=\omega(N_p)\,\kappa_{\mathrm{geom}}\) no longer follows formally.  The
open problem left to \cite{companion15curved} is whether the lower slab/cap
constants can be made uniform near multiply-tangent normal directions, or
whether one must keep constants depending on the fixed curved configuration.

\subsection{Summary of the completed program}

Combining the present paper with its companions, for a $C^2$ strictly convex
body cut by half-spaces in general position the critical scaling of the uniform
doubling condition is, at every transversal stratum of dimension $j$,
\[
\ell_N\;\asymp\;\Bigl(\frac{\log N}{N}\Bigr)^{2/(2n-j)}
\quad(1\le j\le n-1),
\qquad
\ell_N\;\asymp\;\Bigl(\frac{\log N}{N}\Bigr)^{1/n}
\quad(j=0),
\]
exponents and logarithms exact, with constants depending only on the dimension,
the number of cuts, the dihedral-angle lower bound $\theta_0$, and the
curvature band.  The orthogonal model of \cite{companion11} is recovered as the
right-angled, round, flat case $\omega(N_p)=2^{-k}$, $Q=\mathrm{Id}$,
$f_i=e_i$.  The finest resolution at which i.i.d.\ data can certify the
doubling condition is again set by the most isolated extremal sample point of
the worst boundary stratum of the activation geometry, and the explicit
exponent ladder $\tfrac2{n+1},\tfrac2{n+2},\dots,\tfrac1n$ governs the
transition from smooth to conical truncation regardless of the angles at which
the cuts meet.  For curved truncating surfaces, \cite{companion15curved}
settles the fixed-family exponent law and leaves the curvature-band-uniform
lower/slab and factorized-constant problems isolated above.

\subsection*{Acknowledgments}
The research presented in this paper was supported by the European Research
Council (ERC) under the European Union's Horizon Europe research and innovation
programme (grant agreement No.\ 101041711), by the Israel Science Foundation
(grant number 2258/19) and by the Israel Science Foundation (ISF Grant
4101/25).

\subsection*{Declaration on generative AI}
During the preparation of this work the author used generative AI tools to
accelerate drafting and revision.  All mathematical claims, proofs, and
bibliographic information were subsequently reviewed and edited by the author,
who takes full responsibility for the content of the manuscript.

\begin{thebibliography}{9}

\bibitem{ShmaloUDC}
Y.~Shmalo,
\emph{Stability of accuracy for the training of DNNs via the uniform doubling
condition}, Ann.\ Math.\ Artif.\ Intell.\ (2024),
DOI 10.1007/s10472-023-09919-1.

\bibitem{companion2}
Y.~Shmalo, \emph{The uniform doubling condition holds with high probability
for independent samples}, manuscript, 2026.

\bibitem{companion3}
Y.~Shmalo, \emph{The critical scaling of the uniform doubling condition:
smooth bodies versus truncations}, manuscript, 2026.

\bibitem{companion5}
Y.~Shmalo, \emph{The exact critical scaling of the uniform doubling
condition with truncations}, manuscript, 2026.

\bibitem{companion6}
Y.~Shmalo, \emph{The edge exponent of the uniform doubling condition and an
interpolation conjecture}, manuscript, 2026.

\bibitem{companion7}
Y.~Shmalo, \emph{The edge exponent of the uniform doubling condition, II:
the matching upper bound}, manuscript, 2026.

\bibitem{companion11}
Y.~Shmalo, \emph{The interpolation law for the critical scaling of the uniform
doubling condition}, manuscript, 2026.

\bibitem{companion15curved}
Y.~Shmalo, \emph{The curved-truncation companion to the UDC interpolation law},
manuscript, 2026.

\end{thebibliography}

\end{document}
